
\cleardoublepage
\appendix
% --- Start Anhang (A, B, C, ...) ---
\setcounter{chapter}{0} % 1 -> A
\renewcommand{\thechapter}{\Alph{chapter}}
\renewcommand{\thesection}{\thechapter.\arabic{section}}
\renewcommand{\thesubsection}{\thesection.\arabic{subsection}}
\renewcommand{\thesubsubsection}{\thesubsection.\arabic{subsubsection}}

% --- Hyperref: eindeutige Link-Anker (wichtig für korrektes Inhaltsverzeichnis) ---
\renewcommand{\theHchapter}{app.\Alph{chapter}}
\renewcommand{\theHsection}{\theHchapter.\arabic{section}}
\renewcommand{\theHsubsection}{\theHsection.\arabic{subsection}}
\renewcommand{\theHsubsubsection}{\theHsubsection.\arabic{subsubsection}}
% =====================================================
% Anhang A: Mathematische Herleitungen
% =====================================================
\chapter{Mathematische Herleitungen}
\label{app:mathematische-herleitungen}

In diesem Anhang werden die im Haupttext angesprochenen physikalischen Konzepte
formal und mathematisch vertieft. Ziel ist es, die didaktische Lesbarkeit der
Kapitel nicht zu beeinträchtigen und zugleich interessierten Leserinnen und Lesern
die wichtigsten Herleitungen zugänglich zu machen.

Die Abschnitte sind thematisch nach zentralen Eigenschaften und Experimenten des
Elektrons\index{Elektron} gegliedert. Dazu gehören unter anderem das
Ladung-Masse-Verhältnis\index{Ladung-Masse-Verhältnis}, die Bestimmung der
Elementarladung\index{Elementarladung}, die Bewegung geladener Teilchen in
elektrischen und magnetischen Feldern\index{elektrisches Feld}\index{Magnetfeld},
sowie quantenmechanische Eigenschaften wie Spin\index{Spin}, Wellenfunktion\index{Wellenfunktion}
und Aufenthaltswahrscheinlichkeit\index{Aufenthaltswahrscheinlichkeit}.

Auf diese Weise bildet der Anhang eine Brücke zwischen den anschaulichen Erklärungen
im Haupttext und der mathematischen Beschreibung des Elektrons in klassischer Physik,
Quantenmechanik\index{Quantenmechanik}, Relativitätstheorie\index{Relativitätstheorie}
und moderner Teilchenphysik\index{Teilchenphysik}.



%\section{Kapitel I}
\phantomsection
